% !TEX root =  main.tex
%%%%%%%%%%%%%%%%%%%%%%%%%%%%%%%%%%%%%
%%%%%%%%%%%%%%%%%%%%%%%%%%%%%%%%%%%%%
\vspace{-0.2cm}
\section{Related Work}

\textbf{Automated Red-teaming.} Most existing automated red-teaming methods use two forms of feedback signal. One line of research relies mainly on textual feedback or guidance for optimization. These methods benefit from the rich textual information and reasoning capabilities of language models, such as TAP \citep{mehrotra2024tree}, PAP \citep{zeng2024johnny} and AutoDAN-turbo \citep{liu2025autodanturbolifelongagentstrategy}. However, they might suffer from low attack success rate due to a lack of fine-grained quantitative information.

An increasing number of works utilize numerical feedback signals for optimization. Reinforcement learning emerges as a powerful tool to directly optimize a jailbreak model, proven effective by Jailbreak-R1 \citep{guo2025jailbreak} and RL-Hammer \citep{wen2025rl}. However, RL-based approaches face two fundamental challenges: when targeting strong, safety-aligned models, rewards are often sparse, resulting in limited learning signal; when trained against weaker models, learned policies tend to overfit and exhibit poor transferability to more robust targets \citep{liu2025superrlreinforcementlearningsupervision, wen2025rl}. Several methods couple LLM's reasoning capabilities with numerical feedback signals to drive the search or learning process, including PAIR \cite{chao2025jailbreaking} and tree-structured search \citep{sabbaghi2025adversarial}. %Adversarial Reasoning \citep{sabbaghi2025adversarial} explicitly uses LLM reasoning to interpret and act on numerical feedback by utilizing a branching-and-pruning-based search algorithm.
%Additionally, methods such as GCG \citep{zou2023universal} and AutoPrompt \citep{shin2020autoprompt} use logits from target model as signals to guide their prompt optimization. 
Despite employing different optimization signals and algorithms, these red-teaming methods share a common structure: they leverage LLMs' reasoning capabilities through carefully crafted instructions, effectively creating distinct types of agentic systems with manually-designed workflows. However, manually designing the agentic workflow is expensive, time-consuming, and subject to human bias. We explore optimization of red-teaming algorithms from the lens of system design, striving to explore the agentic design space autonomously and efficiently.

%However, these methods often require white-box access to the target model, making them less effective when transfer to black-box settings. Moreover, existing work has manually discovered numerous effective automated red-teaming systems, but it’s challenging to exhaust the ever-evolving target models and benchmarks, further highlighting the importance of automated red-teaming system generation and optimization. Our method utilizes an archive of existing baseline and its pretrained knowledge about red-teaming to actively generate new red-teaming systems. Instead of being a static system, we introduce a pipeline that can learn from the baseline techniques and past generations systems and search for the next generation of red-teaming systems in a dynamic and adaptive way.


\textbf{Agentic Systems Optimization.} Agentic systems aim to solve complicated, long-horizon tasks with structured frameworks, such as self-reflection \citep{madaan2023selfrefineiterativerefinementselffeedback}, multi-persona debate \citep{du2023improvingfactualityreasoninglanguage} and role assignment \citep{xu2025expertpromptinginstructinglargelanguage}. %Agentic systems is a rapidlly expanding research field, 
Their growing use in domains from climate science and healthcare to scientific discovery has motivated work on automating their design, including prompt optimization, hyperparameter tuning, and, most notably, workflow optimization \citep{hou2025investesgmultiagentreinforcementlearning,wang2025medagentproevidencebasedmultimodalmedical,narayanan2024aviary}. Among them, automated workflow optimization aims to optimize the entire workflow structure, offering greater potential because it searches over the full system structure rather than isolated components. Early work such as \citep{hu2025automateddesignagenticsystems} represented workflow as code structure and employed an iterative search algorithm with a linear list to represent its experiences. Later methods \citep{zhang2024aflow} improved search efficiency by introducing a modular approach using a graph-based representation of agentic workflow, utilizing Monte-Carlo Tree Search algorithm to enhance search efficiency \citep{wang2025scoreflow}, or incorporating direct preference optimization \citep{rafailov2024directpreferenceoptimizationlanguage}. These approaches have been proven effective for reasoning tasks such as mathematics and coding \citep{cobbe2021gsm8k,austin2021program}. We extend this paradigm to red-teaming, exploring how automated agentic system optimization can address this unique challenge.

Recent work has begun to apply evolutionary methods to AI safety. EvoSynth targets multi-turn attacks and evolves query-specific attack programs \citep{chen2025evolve}. Claudini applies AutoResearch to red teaming, but targets GCG-style attacks by optimizing discrete suffixes to minimize token-forcing loss \citep{panfilov2026claudiniautoresearchdiscoversstateoftheart,karpathy2026autoresearch}. By contrast, AgenticRed focuses on query-agnostic attack workflows for single-turn red-teaming to produce systems that  generalize across benchmark sets with more diverse attack styles.